\documentclass[12pt, letterpaper]{article}

\usepackage[margin=1in]{geometry}
\usepackage{setspace}
\usepackage{parskip}
\usepackage{enumitem}
\usepackage{titlesec}
\usepackage{hyperref}
\usepackage{microtype}
\usepackage{lmodern}
\usepackage[T1]{fontenc}

\hypersetup{colorlinks=true, linkcolor=blue, urlcolor=blue}

\titleformat{\section}{\large\bfseries}{}{0em}{}
\titlespacing*{\section}{0pt}{1.2em}{0.3em}

\setstretch{1.1}
\setlength{\parskip}{0.5em}

\begin{document}

\begin{center}
    {\LARGE\bfseries Innovation Experiment --- Significance Statement}\\
    \vspace{0.3em}
    {\normalsize Shrinithi Natarajan}
\end{center}

\vspace{0.5em}
\hrule
\vspace{0.8em}

% ============================================================
\section{NIH \#1: Importance of the Problem}

Drug resistance is implicated in the overwhelming majority of cancer deaths in patients with
metastatic disease \cite{liu2025multiomics, mao2025ai}. Despite targeted therapies designed to block
specific oncogenic drivers, acquired resistance emerges in most patients within one to two years of
treatment initiation \cite{mao2025ai}. The problem is not a shortage of data --- public repositories
(GDSC2, CCLE, CTRPv2) contain pharmacological sensitivity profiles matched with genomic,
transcriptomic, and pathway annotations for thousands of cell line--drug pairs, and recent work
demonstrates that learnable general principles of drug sensitivity exist within them
\cite{carli2025cellhit}. What is missing is a framework that reasons over those data layers the
way a trained oncologist does: by recognising that evidence has a natural causal ordering, and that
when layers disagree, the hierarchy --- not a statistical average --- should decide the outcome. No
current computational tool encodes this logic explicitly.

\hrule
\vspace{0.5em}

% ============================================================
\section{NIH \#2: Scientific Premise}

\textbf{We hypothesise that a decentralised multi-agent system resolving inter-modality conflicts
through an explicit five-tier biological Hierarchy of Truth will predict drug resistance more
accurately and more interpretably than monolithic LLM baselines and naive majority-vote systems.}
Three lines of evidence support this.

\textit{First}, the causal hierarchy is already established by biology. Structural DNA alterations
are upstream of transcriptional state, which is upstream of protein activity, which determines
pathway flux \cite{liu2025multiomics, carli2025cellhit}. Encoding this as executable logic --- rather
than leaving it implicit in model weights --- is the core novelty.

\textit{Second}, our preliminary system is working. Four specialised agents (Genomics,
Transcriptomics, Pharmacology, Pathway) each query a curated SQLite database via the Model Context
Protocol \cite{flotho2026mcpmed, widjaja2025bioinformcp}, produce a structured \texttt{EvidencePack},
and feed into a \texttt{DebateEngine} that runs up to three conflict-resolution rounds using a
\texttt{ConflictDetector} and \texttt{AxiomResolver}. All software components are validated by 83 unit and integration tests covering the ETL pipelines,
MCP servers, agent parsing, and protocol logic. The framework will be evaluated against 40
GDSC2-verified cell line--drug pairs once connected to a live LLM. The five tiers are:

\begin{itemize}[leftmargin=1.5em, itemsep=0pt, topsep=2pt]
    \item \textbf{T1 --- Structural Primacy:} DNA-level loss-of-function overrides all downstream signals.
    \item \textbf{T2 --- Transcriptional Gate:} Silencing (z-score $< -2.0$) renders a gene functionally absent.
    \item \textbf{T3 --- Pathway Bypass:} Active bypass routes confer resistance regardless of target status.
    \item \textbf{T4 --- Pharmacological Prior:} Historical IC\textsubscript{50} is a population baseline, overridable by T1--T3.
    \item \textbf{T5 --- Statistical Consensus:} Population statistics apply when molecular evidence is absent.
\end{itemize}

Every prediction will be accompanied by a structured reasoning trace naming the axiom tier that
resolved the case, which agents were overridden, and which dissents were preserved --- a form of
clinical accountability no existing AI tool provides \cite{evomdt2026, han2024consensus,
lorenzo2026oncotimia}.

\textit{Third}, prior approaches have addressable weaknesses. Classical ML methods produce
black-box predictions with no mechanism to inspect which data layer drove the outcome
\cite{saranya2024drncdr, zhao2023msdrp}. LLM-based systems suffer from context overwhelm when
incommensurable data types are presented simultaneously, and from hallucination --- the generation
of biomedical claims not grounded in input data --- a well-documented clinical risk
\cite{kim2025hallucination, nishisako2025hallucination}. Both failure modes are addressed here:
every claim is anchored to a database query, and every conflict is resolved by logged axiom logic.

\hrule
\vspace{0.5em}

% ============================================================
\section{NIH \#3: Improvement to Scientific Knowledge and Clinical Practice}

\textbf{Aim 1} establishes predictive accuracy on the GDSC2 development set (AUROC, AUPRC,
Cohen's $\kappa$, Spearman $\rho$). \textbf{Aim 2} quantifies reasoning quality: hallucination
rate and biological faithfulness score --- evaluation dimensions absent from all prior work
\cite{kim2025hallucination}. \textbf{Aim 3} runs six ablations (no axiom hierarchy; single agent;
free-text data instead of MCP; random tier order; no debate loop; majority vote) to prove each
architectural choice is independently necessary.

Scientific benefit: a reproducible, LLM-agnostic framework for multi-modal biological reasoning.
Technical benefit: an open-source MCP agent-to-tool protocol adaptable to any multi-modal
biomedical task \cite{mehandru2025bioagents}. Clinical benefit: structured reasoning traces that
clinicians can interrogate in the language of molecular oncology --- not a confidence score, but
an auditable argument.

\hrule
\vspace{0.5em}

% ============================================================
\section{NIH \#4: How the Field Will Change}

If the aims are achieved, the field's optimisation target shifts from AUROC points toward reasoning
quality metrics. The axiom-governed debate protocol is generalisable beyond drug resistance to any
domain requiring multi-modal evidence reconciliation --- rare disease diagnosis, treatment response
prediction, biomarker discovery \cite{mehandru2025bioagents}. The debate trace is a novel clinical
artefact: a structured, citable record of which evidence tier resolved each case, directly enabling
clinician oversight of AI recommendations \cite{evomdt2026, han2024consensus}.

\begin{quote}
\small\textit{As a result, our knowledge of how multi-modal molecular evidence should be integrated
during clinical inference will be advanced by an auditable, biologically faithful reasoning
architecture --- one that makes the reasoning process a first-class scientific object and brings
computational oncology closer to AI-assisted decision support that clinicians can trust.}
\end{quote}

\hrule
\vspace{0.5em}

% ============================================================
\section{Reflection}

The Wrenshall sock example taught me to lead with the hypothesis and preliminary data together,
before describing the plan --- reviewers need to believe in the solution before they care about the
details. It also showed that weaknesses belong at the end of the scientific premise, not in a
separate section: strengths first establish credibility, then weaknesses frame exactly what the
aims will address. Switching from fabricated to literature-sourced references sharpened the
statement considerably --- the CellHit, EvoMDT, and MCPmed papers situate this work precisely in
its field rather than citing generic statistics.

% ============================================================
\section{Feedback}

\textit{[To be completed after peer or supervisor review: who gave feedback, how it was solicited,
and what changes were made in response.]}

\vspace{0.8em}
\hrule
\vspace{0.5em}

% ============================================================
{\small
\begin{thebibliography}{99}
\setlength{\itemsep}{0pt}

\bibitem{liu2025multiomics} Liu X et al. Multi-omics advances in understanding cancer drug resistance. \textit{Biomed Technol}. 2025. doi:10.1016/j.bmt.2025.100119

\bibitem{mao2025ai} Mao Y et al. Emerging AI-driven precision therapies in tumor drug resistance. \textit{Mol Cancer}. 2025. doi:10.1186/s12943-025-02321-x

\bibitem{carli2025cellhit} Carli F et al. Learning and actioning general principles of cancer cell drug sensitivity. \textit{Nat Commun}. 2025. doi:10.1038/s41467-025-56827-5

\bibitem{saranya2024drncdr} Saranya KR, Vimina ER. DRN-CDR: cancer drug response prediction using multi-omics. \textit{Comput Biol Chem}. 2024. doi:10.1016/j.compbiolchem.2024.108175

\bibitem{zhao2023msdrp} Zhao H et al. MSDRP: deep learning model for drug response prediction. \textit{Bioinformatics}. 2023.

\bibitem{evomdt2026} Liu Q et al. EvoMDT: self-evolving multi-agent system for clinical decision-making. \textit{npj Digit Med}. 2026. doi:10.1038/s41746-025-02304-8

\bibitem{han2024consensus} Han X et al. Multi-Agent Medical Decision Consensus Matrix for Oncology MDT. \textit{arXiv}:2512.14321. 2024.

\bibitem{lorenzo2026oncotimia} Lorenzo L et al. Evaluation of Oncotimia: LLM system for tumour boards. \textit{arXiv}:2601.19899. 2026.

\bibitem{kim2025hallucination} Kim Y et al. Medical Hallucinations in Foundation Models. \textit{arXiv}:2503.05777. 2025.

\bibitem{nishisako2025hallucination} Nishisako S et al. Reducing Hallucinations in Generative AI for Cancer Information. \textit{JMIR Cancer}. 2025. doi:10.2196/70176

\bibitem{flotho2026mcpmed} Flotho M et al. MCPmed: MCP-enabled bioinformatics for LLM-driven discovery. \textit{Brief Bioinform}. 2026. doi:10.1093/bib/bbag076

\bibitem{widjaja2025bioinformcp} Widjaja F et al. BioinfoMCP: MCP Interfaces in Agentic Bioinformatics. \textit{arXiv}:2510.02139. 2025.

\bibitem{mehandru2025bioagents} Mehandru N et al. BioAgents: multi-agent systems for bioinformatics. \textit{Sci Rep}. 2025. doi:10.1038/s41598-025-25919-z

\end{thebibliography}
}

\end{document}
